\section{Reproducibility}
\label{sec:eval-reproducibility-metric}

The task description names \emph{reproducibility} as one of the evaluation
metrics. Determinism is configured at decoding time
(\texttt{temperature=0}, \texttt{seed=42}) and at the dependency layer
(version-pinned \texttt{package-lock.json}, containerised execution per
Section~\ref{sec:eval-container}); this section reports the
\emph{measured} reproducibility of detector outputs under that
configuration.

\subsection{Method}

A dedicated harness (\texttt{evaluation/reproducibility.js}) takes a
fixed-size subset of the curated corpus (default 10 cases), runs the
analyser $N=3$ times per case with the same seed, and computes two rates:

\begin{itemize}
  \item \textbf{byteIdentityRate} --- the fraction of cases for which every
        run produced byte-identical raw model output. This is the strict
        guarantee deterministic decoding promises.
  \item \textbf{setIdentityRate} --- the fraction of cases for which every
        run produced an identical canonical issue set, i.e.\ identical
        \texttt{(canonical-category, line)} tuples after the canonical
        taxonomy normaliser is applied. This is the more meaningful
        ``the detector's conclusions are reproducible'' metric: it is
        insensitive to harmless paraphrase of issue messages and to
        permutations of the issue array.
\end{itemize}

A pass on \texttt{setIdentityRate} but not on \texttt{byteIdentityRate}
indicates the detector reliably reaches the same conclusion through
slightly different surface text. A failure on
\texttt{setIdentityRate} indicates a non-determinism source that is not
captured by the seed (e.g.\ floating-point scheduling on the GPU).

\subsection{Results}

For \texttt{qwen3:8b} on the locked $N=3$ identical-seed configuration
(10 cases, seed 42, host \texttt{127.0.0.1:11434}), both rates reach
their maximum:

\begin{itemize}
  \item \textbf{\texttt{byteIdentityRate} = 100\,\%} (10 / 10 cases) ---
        every one of the 30 inferences (10 cases $\times$ 3 runs) produced
        byte-identical raw model output. Deterministic decoding with
        \texttt{temperature=0} and a fixed seed therefore holds end-to-end
        on the locked hardware profile.
  \item \textbf{\texttt{setIdentityRate} = 100\,\%} (10 / 10 cases) ---
        the canonical \texttt{(category, line)} issue set is identical
        across all three runs for every case. The detector's conclusions
        are reproducible at the strongest available bar.
\end{itemize}

Median per-inference latency on this sample was 11\,238\,ms; this is
slightly higher than the headline \texttt{qwen3:8b+RAG} median
(1\,573\,ms, Section~\ref{sec:eval-r4}) because the reproducibility
runner uses a more verbose system prompt to be self-contained. Latency
is incidental to the reproducibility metric --- the rate that matters
is whether identical inputs produce identical outputs, which they do.

The byte-identity result is stronger than the literature typically reports
for local LLM inference, where floating-point scheduling on the GPU is
expected to introduce minor non-determinism. The 100\,\% byte-identity
result on Apple M4 Max suggests Ollama's CPU/Metal kernel scheduling is
deterministic under \texttt{temperature=0} on this hardware. Re-runs on
different hardware would re-execute the same harness and produce a
comparable report; we report on the configuration this thesis was
evaluated on, not as a hardware-independent constant.
